\documentclass[10pt]{article}

\usepackage[utf8]{inputenc}

\usepackage[T1]{fontenc}

\usepackage{amsmath}

\usepackage{amsfonts}

\usepackage{amssymb}

\usepackage[version=4]{mhchem}

\usepackage{stmaryrd}

\usepackage{graphicx}

\usepackage[export]{adjustbox}

\graphicspath{ {./images/} }



\begin{document}

становится практически неосуществимым, причем использование самых быстрых вычислительных машин лишь несущественно распиряет границы его применимости. Это означает, что требуется дальнейшее уточнение постановки C. 3.



Сложившиеся здесь тенденции продемонстрированы ниже на примере одного из классов у. с. - схем из gанкциональных элементов для одного частного случая, когда сложность понимается как число әлементов в схеме. Функции, реализуемые этим классом у. с.,это функции алгебры логики (булевы функции).



Прежде всего следует отметить, что неприемлемость решения, к-рое дает тривиальный метод,- это лишь одна из ігричин, ведущих к пересмотру постановки С. 3. Существует гипотеза, что любой алгоритм, к-рый для каждої булевой функции строит схему с минимальным числом элементов, неизбежно содержит в себе в каком-то смысле перебор всех булевых функций. В пользу этой гипотезы говорит теорема о том, что если рассмотреть бесконечное множество булевых функций, содержащее при каждом числе переменных по одной самой сложной функции, и замкнуть это множество относительно операций переименования переменных (без отождествления) и подстановки констант, то получится множество, состоящее из всех вообще булевых функций (с точностью до несущественных переменных функции) Эта гипотеза получила много других косвенных подтверждений, причем нек-рые из них свидетельствуют о том, что если даже ослабить требование по числу әлементов, допуская построение схем, близких к оптимальным, но сохранить требование применимости алгоритма к каждой булевой функции, то перебор не должен существенно уменьшиться. Идея о принципиальной неизбежности перебора и является основной причиной видоизменения постановки С. 3. Здесь имеется несколько возможностей.



Шенноновский подход (по имени $\nprec$. Э. Шеннона, к-рыї впервые предложил его для контактных схем) связан с рассмотрением функции



\[

L(n)=\max L(f)

\]



где $L(f)$ - минимально возможное число элементов в схеме, реализующей функцию $f$, а максимум берется по всем булевым функциям $f$ от $n$ переменных. При этом $L(f)$ и $L(n)$ зависят от базиса, т. е. от множества функциональных элементов, к-рые можно использовать при построении схем. С. з. состоит в том, чтобы найти эффективный метод построения для каждой булевой функции от $n$ переменных схем с числом әлементов, существенно не превосходящим $L(n)$.



При более общей постановке задачи каждому функциональному элементу $E$ сопоставляется положительное число $L(E)$ - сложность этого элемента, к-рая зависит только от функции, реализуемой элементом $E$, а каждой схеме $S$ приписывается сложность



\[

L(S)=\sum L(E)

\]



где суммирование производится по всем әлементам схемы $S$. Для каждой булевой функции $f$ определяется сложность ее реализации (в заданном базисе):



\[

L(f)=\min L(S)

\]



где минимум берется по всем схемам $S$, реализующим функцию $f$. Затем, как и раньше, вводится $L(n)=$ $=\max L(f)$ и требуется для произвольной булевой функции от $n$ переменных построить схему, имеюгцу сложность не выше (или не на много выше), чем $L(n)$.



При таком подходе нужно уметь вычислять или хорошо оценивать хотя бы снизу величину $L(n)$. Оказывается, что неплохую нижнюю оценку для $L(n)$ дает уже м о ш о с т ш о й м е т о д, основанный на следующем соображении: число всевозможных схем, реализующих булевы функции от $n$ переменных и имегщих сложность не более $L(n)$, должно быть не меньше числа всех булевых функций от $n$ псременных. В мопностном методе требуется только получить удовлетворительную верхнюю оценку для числа рассматриваемых схем.



После этого о качестве любого универсального метода синтеза (т. е. метода, пригодного для Іроизвольной булевой функции) можно в значительной мере судить по тому, во сколько раз сложность построенных им схем при соответствующем $n$ отличается от нижней оценки для $L(n)$. Одновременно сложность әтих схем служит конструктивной верхней оценкой для $L(n)$. Основным достижением в этом направлении является создание эффективного универсального метода синтеза, к-рый дает схемы со сложностью, мало отличающейся от нижней оценки для $L(n)$, а при $n \rightarrow \infty$ асимптотически совпадающей с ней. Таким образом, это асимптотически наилучший метод синтеза, позволяющий к тому же установить асимптотич. поведение функции $L(n)$. Пусть



\[

\rho=\min L(E) /(i-1),

\]



где $i$ - число входов әлемента $E$, а минимум берется по всем әлементам, имеющим не менее двух входов. Тогда



\[

L(n) \sim \rho \frac{2^{n}}{n} .

\]



Для других классов у. с. слоякность $L(S)$, а с ней и сложности $L(f)$ и $L(n)$ вводятся аналогичным образом. Обычно $L(S)$ для контактной схемы $S$ есть число контактов в ней, а для формулы $S$ - это число символов переменных в ней. Для автоматов и машин Тьюринга $L(S)$ иногда вводится как число состояний автомата $S$ или машины Тьюринга $S$. Соответствующее асимптотическое поведение величины $L(n)$ для этих классов показано в таблице 1.



Т а б лиц а 1



\begin{center}

\includegraphics[max width=\textwidth]{2023_07_08_282b2ee3ae571cdc4182g-1}

\end{center}



Если понятие сложности вводится иначе, напр. как вероятность ошибки, или если рассматривается другой класс у. с., напр. автоматы, то соответствуюшее асимптотич. соотношение может содержать не один, а дьа и более параметров, причем в общем случае задача о получении асимптотики, напр. для автоматов, алгоритмически неразрешима.



Анализ нижней оценки для $L(n)$ показывает, что на самом деле для почти всех булевых функций от $n$ переменных сложность реализации асимптотически не меньше $L(n)$. Это означает, что данный метод синтеза для почти всех булевых функций дает почти самые простые схемы.





\end{document}